\documentclass[11pt,a4paper]{article}
\usepackage[margin=25mm]{geometry}
\usepackage{amsmath,amssymb,amsthm,bm,mathtools}
\usepackage[T1]{fontenc}
\usepackage{lmodern,microtype}
\usepackage[colorlinks=true,linkcolor=blue,urlcolor=blue]{hyperref}
\newcommand{\D}{\mathcal D}
\newcommand{\W}{\mathcal W}
\newcommand{\Sact}{\mathcal S}
\newcommand{\E}{\mathbb E}
\title{Section 1 note: the conjugate (backward) saddle-point equation}
\author{}
\date{}
\begin{document}
\maketitle

\begin{abstract}
For the hidden layers, Section~1 gives the conjugate saddle point
(\texttt{eq:saddle-phi-hat}) as
$\hat\Phi^{(l)} = \frac1N\langle\hat h\hat h\rangle
= \frac1N\langle\partial_\eta h\,\partial_\eta h\rangle$
and, equivalently, $\hat\Phi^{(l)} = \Phi^{-1}\langle hh\rangle\Phi^{-1}$
(\texttt{eq:backward-pass}). Both are wrong. Integrating out $\hat h$ leaves a
$\Phi$-dependent normalisation $\det(2\pi\Phi)^{-1/2}$ that was dropped; keeping
it adds $-\Phi^{-1}$ to the backward pass. The response correlator enters with a
minus sign, and the ``derivative'' interpretation needs a curvature term. The
correct equation is
\begin{align*}
  \hat\Phi^{(l)} &= -\frac1N\langle\hat h\hat h\rangle^{(l+1)}
  = \frac1N\Big(\Phi^{-1}\langle h h^\top\rangle^{(l+1)}\Phi^{-1} - \Phi^{-1}\Big)\\
  &= \frac1N\big\langle\nabla V\,\nabla V^\top + \nabla\nabla^\top V\big\rangle^{(l+1)},
\end{align*}
with $V$ the tilt acting on layer $l+1$. As written, the backward pass gives a
nonzero conjugate even when no signal arrives from downstream. The one-hidden-layer
result \texttt{eq:linear-soln} does not use this equation and is unaffected;
deep networks and the effective-network picture of Appendix~A1 are affected.
\end{abstract}

\section{Setting}

Use the rotated (real) convention of \texttt{setup\_conjugate\_rotation.tex}
and the joint layer average of \texttt{setup\_site\_factorisation.tex}. For
$1\le l\le L-1$, $\Phi\equiv\Phi^{(l)}$ enters the action only through the
$\frac12\hat\Phi^{(l)}\cdot\Phi^{(l)}$ term and through the single site of layer
$l+1$, for which it is the prior covariance:
\begin{align}
  \Sact &\ni \tfrac12\hat\Phi^{(l)}\cdot\Phi - \tfrac1N\log\W[\Phi,\hat\Phi^{(l+1)}],\\
  \W[\Phi,\hat\Phi^{(l+1)}]
  &= \int\D h\,\D\hat h\,\exp\Big\{-i\hat h\cdot h - \tfrac12\hat h\Phi\hat h + V(h)\Big\},
  \label{eq:W-msr}\\
  V(h) &= \tfrac12\phi(h)^\top\big(\hat\Phi^{(l+1)}K^{(l+1)} - G^{(l+1)}\big)\phi(h),
\end{align}
where $h$ collects all $NP$ preactivations of layer $l+1$ and the contraction of
$\hat\Phi^{(l+1)}K^{(l+1)}$ is over the post-synaptic slots, as in the single site.

\section{The dropped normalisation}

The $\hat h$ integral in \eqref{eq:W-msr} is Gaussian:
\begin{equation}
  \int\D\hat h\,e^{-i\hat h\cdot h - \frac12\hat h\Phi\hat h}
  = \det(2\pi\Phi)^{-1/2}\,e^{-\frac12 h\Phi^{-1}h}
  \quad(\text{with }\D\hat h = \textstyle\prod d\hat h/2\pi),
\end{equation}
so
\begin{equation}
  \W = \det(2\pi\Phi)^{-1/2}\int\D h\,e^{-\frac12h\Phi^{-1}h + V(h)}.
  \label{eq:W-h}
\end{equation}
Section~1 writes the second line of \texttt{eq:single-site} without the
prefactor. The prefactor is irrelevant for averages \emph{at fixed} $\Phi$ (so
the forward equation for $\Phi^{(l)}$ is unaffected), but $\Phi$ is now a
variable of the action and the prefactor contributes
$-\frac12\log\det\Phi$ to $\log\W$. It is $O(N)$, the same order as the rest of
$\log\W$, so it cannot be dropped at the saddle point.

\section{Derivation of the correct equation}

\paragraph{From the $h$ form.} Differentiating \eqref{eq:W-h} (treating
$\Phi$ as symmetric; the same convention on both sides),
\begin{equation}
  \frac{\partial\log\W}{\partial\Phi}
  = -\tfrac12\Phi^{-1} + \tfrac12\Phi^{-1}\langle hh^\top\rangle\Phi^{-1}.
\end{equation}
Then $0 = \partial\Sact/\partial\Phi^{(l)} = \frac12\hat\Phi^{(l)} - \frac1N\partial_\Phi\log\W$ gives
\begin{equation}
  \boxed{\;\hat\Phi^{(l)} = \frac1N\Big([\Phi^{(l)}]^{-1}\langle hh^\top\rangle^{(l+1)}[\Phi^{(l)}]^{-1} - [\Phi^{(l)}]^{-1}\Big).\;}
  \label{eq:bp}
\end{equation}
The $-\Phi^{-1}$ comes entirely from the prefactor. Section~1's form,
$\Phi^{-1}\langle hh\rangle\Phi^{-1}$, lacks it and also lacks the overall $1/N$
(it had $\frac1N\sum_k\langle\cdot\rangle_k$ instead; see the factorisation note).

\paragraph{From the $\hat h$ form.} Differentiating \eqref{eq:W-msr} directly,
$\partial_\Phi\log\W = -\frac12\langle\hat h\hat h^\top\rangle$, hence
\begin{equation}
  \hat\Phi^{(l)} = -\frac1N\langle\hat h\hat h^\top\rangle^{(l+1)},
  \qquad
  -\langle\hat h\hat h^\top\rangle = \Phi^{-1}\langle hh^\top\rangle\Phi^{-1} - \Phi^{-1}.
  \label{eq:hh}
\end{equation}
Section~1 has $+\frac1N\langle\hat h\hat h\rangle$. The sign matters: the MSR
response correlator $\langle\hat h\hat h\rangle$ vanishes for an untilted
Gaussian and is negative semi-definite whenever the tilt broadens $h$, while the
readout conjugate $\hat\Phi^{(L)} = \tilde\Phi^{-1}\mathbb Y\tilde\Phi^{-1}$ is
positive semi-definite; only the minus sign makes the hidden and readout
conjugates the same kind of object. (In the unrotated MSR convention the
equation reads $i\hat\Phi^{(l)} = -\frac1N\langle\hat h\hat h\rangle$, so the
section's form is wrong in either convention.)

\paragraph{The ``derivative'' interpretation.}
Section~1 identifies $\langle\hat h\hat h\rangle$ with
$\langle\partial_\eta h\,\partial_\eta h\rangle$, $\eta$ the input current. For an
additive current $\partial_\eta h = 1$, so this cannot be meant literally. The
precise statement is Gaussian integration by parts for the tilted measure
$p(h)\propto\mathcal N(h;0,\Phi)\,e^{V(h)}$: from
$\E_p[h\,f(h)] = \Phi\,\E_p[\nabla f + f\,\nabla V]$ applied to $f = h^\top$ and then to
$f = \nabla V^\top$,
\begin{equation}
  \E_p[hh^\top] = \Phi + \Phi\,\E_p\big[\nabla V\nabla V^\top + \nabla\nabla^\top V\big]\Phi,
\end{equation}
and so, combining with \eqref{eq:bp},
\begin{equation}
  \hat\Phi^{(l)\mu\nu}_{ij} = \frac1N\Big\langle
    \frac{\partial V}{\partial h^\mu_i}\frac{\partial V}{\partial h^\nu_j}
    + \frac{\partial^2V}{\partial h^\mu_i\,\partial h^\nu_j}\Big\rangle^{(l+1)} .
  \label{eq:stein}
\end{equation}
$\partial V/\partial h$ is the force that the downstream constraint (the tilt from
layer $l+2$, or from the readout) exerts on a preactivation of layer $l+1$: the
analogue of a back-propagated error, which is presumably what was meant. The
second term, the mean curvature of the tilt, has no counterpart in the
section's formula and is not small in general: for
$V = \frac12 b\,h^2$ (linear $\phi$, one scalar site, $\Phi = 1$) the first term
gives $b^2/(1-b)$ and the second $b$, summing to the exact $b/(1-b)$.

\paragraph{Numerical check.} For two coupled sites with
$\Phi = \begin{psmallmatrix}1&0.4\\0.4&0.8\end{psmallmatrix}$, $\phi=\tanh$ and
$V = \frac12\phi^\top B\phi$ with
$B = \begin{psmallmatrix}0.6&-0.2\\-0.2&0.3\end{psmallmatrix}$, grid quadrature gives
$\partial_\Phi\log\W$ (by finite differences of \eqref{eq:W-h}) equal to
$\frac12(\Phi^{-1}\langle hh^\top\rangle\Phi^{-1}-\Phi^{-1})$ to $3\times10^{-11}$,
and the latter equal to the right-hand side of \eqref{eq:stein} to $5\times10^{-13}$.
Without the determinant, the finite difference misses the $-\frac12\Phi^{-1}$ term.

\section{Consequences}

\begin{itemize}
  \item \textbf{Untilted limit.} If layer $l+1$ receives no signal from downstream
    ($V\equiv0$), then $\langle hh^\top\rangle = \Phi$ and \eqref{eq:bp} gives
    $\hat\Phi^{(l)} = 0$, as it must. The section's formula gives
    $\hat\Phi^{(l)} = \Phi^{-1}$: a spurious tilt, of order one, that would
    propagate down every layer of a deep network.
  \item \textbf{Deep networks.} The hidden-layer saddle point
    ($L\ge2$) changes. Any numerical solver of the forward/backward
    equations must use \eqref{eq:bp}.
  \item \textbf{Section~2 / \texttt{eq:linear-soln}.} Not affected: with one
    hidden layer the only conjugate is $\hat\Phi^{(1)}=\hat\Phi^{(L)}$, fixed by the
    readout, and the forward equation uses normalised averages, which do not see
    the prefactor.
  \item \textbf{Appendix A1.} The ``backward pass'' of the effective network
    cites \texttt{eq:backward-pass}, and two interpretive remarks (that $\hat\Phi$
    is the correlation of derivatives with respect to inputs, and that $B$
    encodes the covariance of preactivations with respect to their inputs) cite
    \texttt{eq:saddle-phi-hat}. Both should be read with \eqref{eq:bp} and
    \eqref{eq:stein}: the conjugate is the mean outer product of downstream
    forces plus their mean curvature.
\end{itemize}

\section{Recommended fix}

In Section~1 (done, in red): restore $\det(2\pi\Phi)^{-1/2}$ in the second line of
\texttt{eq:single-site} (and in the product form of $\W$); replace
\texttt{eq:saddle-phi-hat} by \eqref{eq:hh}--\eqref{eq:stein}; replace
\texttt{eq:backward-pass} by \eqref{eq:bp}.

\end{document}
